\title{Probing Model Quality in an Intelligence Market: Calibrated Canaries in Noesis}

\author{Giacinto Paolo (GP) Saggese$^{1}$}
\author{Javin Ahuja$^{1}$}
% TODO(*): Fill in the affiliation of Giles Greene.
\author{Giles Greene}
\author{Lucas Caicedo$^{2}$}
\thanks{$^{1}$ University of Maryland, College Park, MD 20742, USA.
$^{2}$ The George Washington University, Washington, DC, USA}

\date{\today}

\begin{abstract}
A buyer of machine intelligence cannot observe which model actually served a
request, and a seller can lower its cost by serving a weaker model under the
name of the contracted one. The \Noesis{} protocol addresses this gap with
fulfillment monitoring, which requires a per-request estimate of the capability
a seller delivered. This paper describes the canary subsystem that implements
that estimate in the \Noesis{} prototype. A canary is a known-answer question
that the platform sends to a seller as ordinary buyer traffic and grades
deterministically. The subsystem comprises a versioned bank of 300 generated
candidate questions in six task families; a deterministic grader with
three-valued outcomes; a calibration procedure that keeps a question only if
every healthy provider answers it correctly in every attempt and a weaker model
answers it correctly in at most 60\% of attempts; keyed, family-balanced
selection and scheduling of probes; a prober that persists graded evidence; and
a verdict rule that flags a seller only when a one-sided confidence upper bound
on its canary accuracy falls below the promise. An exact binomial analysis
indicates that, with ten canaries and a promised accuracy of 0.85, a seller
whose true accuracy is 0.5 receives a failing quality verdict with probability
0.83, whereas a seller at 0.95 does so with probability 0.001. At the time of
writing, the live provider calibration has not been executed, so the paper
reports the design, its analytical properties, and its unit-level validation,
and no measured model accuracies.
\end{abstract}

\maketitle

\setcounter{tocdepth}{2}
\tableofcontents

% ##############################################################################
\section{Introduction}
\label{sec:introduction}

Applications and agents buy large language model (LLM) inference as a named
product, such as a specific model served by a specific provider. The buyer
observes the text of a response and its latency, but not the model weights, the
numerical precision, or the serving stack that produced it. Two lines of
evidence suggest that this opacity has consequences. First, the behavior of a
model served under a fixed name has been shown to change over time on
well-defined tasks \cite{chen2023drift}. Second, recent work audits whether
endpoints serve the model they advertise and reports endpoints whose output
distributions differ from the reference weights \cite{gao2025equality}, and
discusses the incentive for providers to substitute cheaper or compressed
models while charging the advertised price \cite{cai2025substitution}.

The \Noesis{} white paper \cite{noesis_whitepaper} proposes to trade machine
intelligence as a commodity through two components. \NoesisMarket{} clears
bids and asks into intelligence contracts. \NoesisServer{} executes those
contracts against real providers and meters whether the delivered capability,
latency, and reliability matched the contract. The metering step needs a
per-request estimate $\hat c(x)$ of the capability tier delivered. The white
paper deliberately leaves the estimator pluggable and names ``a verifier model,
a benchmark probe, or agreement with a stronger reference model'' as candidate
implementations. It also lists the reliability of the per-request estimates
among its open questions, because the reputation mechanism of the market is
only as trustworthy as the monitoring that feeds it.

This paper documents the first concrete instantiation of that estimator, the
\emph{canary subsystem} of the \Noesis{} prototype. The term follows the
engineering usage of a small, controlled probe whose outcome is known in
advance; it is unrelated to the use of canary strings for measuring
memorization in training data \cite{carlini2019secret}. The subsystem lives in
the \texttt{checker/} and \texttt{prober/} directories of the \texttt{Noesis}
research code base, and this paper describes the state of that code at commit
\texttt{e5f730c3}.

\subsection{Contributions}
\label{sec:contributions}

\begin{enumerate}
  \item \textbf{Placement in the protocol.} The paper states how known-answer
    probes instantiate the capability-measurement component of the white
    paper and how their outcomes enter the contract verdict, including the
    separation of the white paper's joint success indicator into a latency
    metric and a quality metric (Section~\ref{sec:background}).

  \item \textbf{Calibrated canaries.} The paper specifies the canary bank, the
    deterministic grader, and a calibration protocol that measures, before use,
    which candidate questions are answered reliably by healthy providers and
    unreliably by a weaker model. It derives closed-form survival
    probabilities for a candidate (Sections~\ref{sec:bank} and
    \ref{sec:calibration}).

  \item \textbf{Hidden, balanced probing.} The paper describes keyed
    selection and scheduling of probes, proves bounds on family balance and on
    probe spacing, and describes how the prober makes probes indistinguishable
    from buyer traffic at the provider boundary while remaining safe under
    concurrent workers and partial failure (Section~\ref{sec:probing}).

  \item \textbf{Verdict statistics.} The paper defines the evidence rules and
    the one-sided upper-bound test, quantifies detection power and false-alarm
    probability by exact enumeration, and compares the rule with the lower-bound
    rule of the white paper (Section~\ref{sec:verdict}).

  \item \textbf{Status and limits.} The paper reports which parts of the
    system are validated by tests, states that the paid live calibration is
    pending, and lists the known limitations of the design
    (Sections~\ref{sec:validation} and \ref{sec:limitations}).
\end{enumerate}

\subsection{Scope}
\label{sec:scope}

The paper covers the production of quality evidence and its use in the quality
half of the verdict. It does not cover the auction, the reputation update, the
gateway's buyer-facing path, or the window-closing worker that freezes evidence
and writes a verdict row, which the code base lists as remaining integration
work. Quantities in tables and in the text come from the source code, from the
design documents in \texttt{research/Noesis/docs}, or from the exact
computations in \texttt{figures/make\_figures.py}; none comes from a live
provider run.

\subsection{Organization}

Section~\ref{sec:background} recalls the white paper's contract and monitoring
notation, and states the problem and the threat model.
Section~\ref{sec:architecture} gives the architecture and the persisted data
model. Section~\ref{sec:bank} describes the bank and the grader.
Section~\ref{sec:calibration} describes calibration.
Section~\ref{sec:probing} describes selection, scheduling, and the prober.
Section~\ref{sec:verdict} turns probe outcomes into verdicts and analyzes
detection power. Section~\ref{sec:validation} reports validation status,
Section~\ref{sec:related_work} discusses related work, and
Section~\ref{sec:limitations} lists limitations and future work.
